\section{Resultados esperados}
Para abordar la sección de resultados esperados, se puede consultar la tabla ~\ref{tab:resultados}, la cual se basa en los objetivos específicos. Al cumplir estos objetivos específicos, se logrará alcanzar el objetivo general.

\begingroup
\setlength{\tabcolsep}{10pt} % padding horizontal
\setlength{\extrarowheight}{3pt} 
\renewcommand{\arraystretch}{1.25}

\begin{longtable}{|p{8cm}|p{8cm}|}
\caption{Resultados esperados según los objetivos específicos}
\label{tab:resultados} \\
\hline
\textbf{Objetivo específico} & \textbf{Resultado esperado} \\
\hline
\endfirsthead

\hline
\textbf{Objetivo específico} & \textbf{Resultado esperado} \\
\hline
\endhead

\endfoot

\hline \hline
\endlastfoot

1. Identificar las teorías y enfoques clave relacionados con el aprendizaje autónomo y continuo, así como las estrategias efectivas para fortalecer la capacidad de aprender a aprender. 
& \begin{itemize}[leftmargin=*, topsep=0pt, partopsep=0pt, parsep=0pt, itemsep=0.4\baselineskip]
     \item Documento sobre los fundamentos teóricos y prácticos que sustenten el desarrollo del módulo “Learning to learn”
    \item Selección y adopción de las estrategias más adecuadas para fortalecer la capacidad de aprender a aprender de los estudiantes de ingeniería de sistemas de la Universidad del Valle, sede Tuluá.
\end{itemize}
\\ \hline

2. Investigar metodologías y técnicas empleadas para fomentar la escucha activa durante las interacciones comunicativas, considerando la literatura académica y las mejores prácticas en el campo de la comunicación interpersonal.
& \begin{itemize}[leftmargin=*, topsep=0pt, partopsep=0pt, parsep=0pt, itemsep=0.4\baselineskip]
     \item Documento sobre las diferentes técnicas y enfoques utilizados para fomentar la escucha activa.
    \item El diseño de actividades y estrategias en el módulo de “Active listening”
\end{itemize}
\\ \hline

3. Desarrollar actividades gamificadas y experiencias de aprendizaje interactivas, utilizando recursos educativos relevantes que respalden el desarrollo de las habilidades blandas de “aprender a aprender” y “escucha activa”. 
& \begin{itemize}[leftmargin=*, topsep=0pt, partopsep=0pt, parsep=0pt, itemsep=0.4\baselineskip]
     \item Conjunto de recursos educativos adaptados y seleccionados específicamente para respaldar el desarrollo de las habilidades blandas de “aprender a aprender” y “escucha activa”
    \item Conjunto de actividades gamificadas y experiencias de aprendizaje interactivas, que incluya instrucciones detalladas para cada actividad, así como materiales necesarios y pautas para la evaluación del aprendizaje y el progreso de los estudiantes. 
\end{itemize}
\\ \hline

4. Implementar los módulos de “aprender a aprender” y “escucha activa”. 
& \begin{itemize}[leftmargin=*, topsep=0pt, partopsep=0pt, parsep=0pt, itemsep=0.4\baselineskip]
    \item Diseños MER de la base de datos: Modelos Entidad-Relación (MER) detallados para la base de datos que soportará la aplicación.
    \item Integración de modelos de procesamiento de lenguaje natural (PLN): Implementación de modelos de OpenAI para la generación de contenidos personalizados, incluyendo roadmaps de aprendizaje adaptados al usuario en el módulo Aprender a Aprender, así como la creación dinámica de textos narrativos y diálogos para los desafíos del módulo de Escucha Activa, los cuales se integran con ElevenLabs para la síntesis y reproducción de audio.
    \item Documentación de la arquitectura de los módulos.
    \item El código fuente de los módulos.
\end{itemize}
\\ \hline

5. Evaluar la efectividad de los módulos en términos de comprensión del contenido, participación, motivación, retención de información y aplicación práctica, a través de pruebas piloto y retroalimentación de los estudiantes, identificando áreas de mejora.
& \begin{itemize}[leftmargin=*, topsep=0pt, partopsep=0pt, parsep=0pt, itemsep=0.4\baselineskip]
    \item Aplicación de encuestas de percepción de mejora a los usuarios que utilicen ambos módulos, con el propósito de evaluar su nivel de satisfacción, motivación y percepción de progreso en habilidades como la escucha activa y el aprendizaje autónomo.
    \item Un informe con el análisis y resultados de las pruebas realizadas, con conclusiones y recomendaciones para futuras mejoras.
\end{itemize}
\\ \hline

\end{longtable}